\section{Reward Granularity and the \texorpdfstring{$K{>}2$}{K>2} Collapse}
\label{appendix:reward_discrimination}

ParetoBandit evaluates a $K{=}3$ portfolio, while some prior systems route across larger pools: PILOT~\cite{panda2025pilot} uses $K{=}11$ on Routerbench and PROTEUS~\cite{bhatti2026proteus} trains on $K{=}11$ (Routerbench) and $K{=}14$ (SPROUT).  We argue that \emph{reward granularity}, not portfolio size, is the binding constraint for genuine multi-model routing, and that $K{>}2$ pools under binary rewards collapse to an effective $K{=}2$ strong-vs.-weak triage.

\paragraph{Collapse mechanism.}
Modern LLMs at adjacent price tiers produce acceptable responses on the majority of prompts.  A binary reward signal (correct/incorrect, win/loss) therefore clusters frontier and mid-tier models into a single ``strong'' tier whenever both exceed the acceptance threshold.  In our $K{=}3$ portfolio, a natural binarisation threshold ($\tau{=}0.85$) causes the two stronger models to tie on $85.8\%$ of prompts, eliminating the router's ability to distinguish quality within a budget.  Switching to the continuous composite reward ($r_t \in [0,1]$) reduced tied prompts to ${<}1\%$ and raised best-arm entropy from ${\sim}1.1$ to ${\sim}2.7$ bits.

\paragraph{Evidence from PILOT ($K{=}11$, exact-match rewards).}
PILOT's qualitative routing analysis (Section~5.1 of~\cite{panda2025pilot}) confirms this collapse at scale.  Despite $K{=}11$ arms, per-task traffic concentrates on a single model: $90\%$ GPT-4 on MMLU, $89.4\%$ GPT-4 on ARC, $94\%$ Claude-v1 on GSM8K, and at most two active models on MBPP.  The remaining 9--10 arms carry negligible traffic.  The appearance of diverse $K{=}11$ routing is an artefact of aggregating across heterogeneous tasks; within each task, exact-match rewards collapse the pool to one or two effective arms---the $K{=}2$ triage our analysis predicts.

\paragraph{Literature pattern.}
The pattern is consistent across the routing literature.  RouteLLM~\cite{ong2024routellm} frames routing as an explicit binary strong/weak decision and acknowledges the $K{=}2$ restriction as a limitation.  LLMRouterBench~\cite{li2026llmrouterbench} identifies binary metrics and dominant-model pools as key evaluation failures for multi-model routing.  Conversely, systems that demonstrate genuine multi-model discrimination universally adopt continuous rewards: BaRP~\cite{wang2025barp} ($K{=}11$ on Routerbench, $K{=}6$ on NQ/HotpotQA, scalarized $r - \lambda c$), MixLLM~\cite{wang2025mixllm} ($K{=}11$, predicted accuracy on Routerbench), and PROTEUS~\cite{bhatti2026proteus} ($K{=}11$ on Routerbench, $K{=}14$ on SPROUT, continuous SLA targets).  Notably, MixLLM evaluates on the same $K{=}11$ Routerbench pool as PILOT, yet its continuous quality predictor avoids the single-model collapse that PILOT exhibits; however, neither MixLLM nor BaRP reports per-model traffic distributions, so the degree to which all eleven arms carry meaningful traffic remains unverified.

\paragraph{Implication.}
The constraint is not portfolio size but reward resolution.  A $K{=}11$ pool with binary rewards produces at most $K{=}2$ effective routing decisions per task; a $K{=}3$ pool with continuous rewards produces genuine three-way routing that varies with prompt context and budget.  Extending ParetoBandit to larger $K$ is straightforward---the LinUCB update is $O(Kd^2)$, and the onboarding experiment (\S\ref{sec:eval_onboarding}) already demonstrates $K{=}3 \to K{=}4$ transitions---and the continuous reward signal ensures additional models will be discriminated rather than collapsed.
